\section{问题三模型的建立与求解}

\subsection{问题三的描述与决策变量}

问题三回到问题一的室外条件：环境温度为$-40\,^{\circ}\mathrm C$、无风且实验者静止不动，因此人体侧对流系数、服装外表面自然对流关系、人体等效热容以及三层瞬态传热方程均继承问题一。新的设计任务是增加防护服厚度，同时满足材料费用和负重能力约束，使人体等效表面温度首次降至$15\,^{\circ}\mathrm C$的时间尽可能长。

附件明确规定：中间功能层最大厚度为$0.45\ \mathrm{mm}$，但“后期增加厚度只调整最外层厚度”，且最外层涂层每层厚度固定为$0.3\ \mathrm{mm}$。因此本问不把三层厚度同时作为连续变量，而只增加离散的最外隔热涂层。令$n$为新增最外层涂层数，$n\in\mathbb Z_{\ge0}$，则
\begin{equation}
d_1=0.7\ \mathrm{mm},\qquad
d_2=0.4\ \mathrm{mm},\qquad
d_3(n)=0.3(1+n)\ \mathrm{mm}.
\label{eq:q3-thickness}
\end{equation}
其中原有$0.3\ \mathrm{mm}$最外层对应$n=0$。总厚度为
\begin{equation}
L(n)=1.4+0.3n\ \mathrm{mm}.
\end{equation}
由于$d_2=0.4\ \mathrm{mm}<0.45\ \mathrm{mm}$且本问不再增加中间层，功能层硬度约束自动满足。

\subsection{材料费用约束}

沿用问题一的有效传热面积$A=A_b=1.6521\ \mathrm{m^2}$。附件给出内层织物价格为$500$元/$500\ \mathrm g$，最外隔热层价格为$150$元/$500\ \mathrm g$，对应单位质量价格
\begin{equation}
p_1=1000\ \mathrm{yuan/kg},\qquad
p_3=300\ \mathrm{yuan/kg}.
\end{equation}
中间功能层价格为$10\ \mathrm{yuan/m^2}$。原防护服三层材料费用为
\begin{equation}
C_0
=A\rho_1d_1p_1+10A+A\rho_3d_3(0)p_3.
\label{eq:q3-base-cost}
\end{equation}
代入附件参数可得
\begin{equation}
C_0\approx301.67\ \mathrm{yuan}.
\end{equation}
每新增一层$0.3\ \mathrm{mm}$最外涂层，增加质量
\begin{equation}
\Delta m=A\rho_3(0.3\times10^{-3})\approx0.148689\ \mathrm{kg},
\label{eq:q3-layer-mass}
\end{equation}
相应增加费用
\begin{equation}
\Delta C=p_3\Delta m\approx44.61\ \mathrm{yuan}.
\label{eq:q3-layer-cost}
\end{equation}
题目允许在原支出基础上增加不超过$50\%$的资金，因此
\begin{equation}
C_0+n\Delta C\le1.5C_0.
\end{equation}
从而
\begin{equation}
n\le
\left\lfloor\frac{0.5C_0}{\Delta C}\right\rfloor
=3.
\label{eq:q3-budget-nmax}
\end{equation}
因此费用约束已经把设计空间压缩为有限集合
\begin{equation}
\boxed{n\in\{0,1,2,3\}},
\end{equation}
对应最外层厚度分别为$0.3,0.6,0.9,1.2\ \mathrm{mm}$。本问无需再使用连续优化或启发式算法，对四个候选方案逐一进行传热正演即可得到全局最优解。

\subsection{随站立时间衰减的负重约束}

三层防护服在方案$n$下的总质量为
\begin{equation}
m_c(n)
=A\left[\rho_1d_1+\rho_2d_2+\rho_3d_3(n)\right].
\label{eq:q3-clothing-mass}
\end{equation}
其中基准服装质量约为
\begin{equation}
m_c(0)\approx0.7542\ \mathrm{kg},
\end{equation}
每新增一层最外涂层增加$\Delta m\approx0.148689\ \mathrm{kg}$。

题目规定实验者初始最大承受重量为$100\ \mathrm{kg}$，且每多承受$10\ \mathrm s$，最大承受重量下降$0.5\ \mathrm{kg}$。将站立时由人体支撑的总质量解释为人体自身与防护服质量之和，则时刻$t$的承重能力为
\begin{equation}
M_{\max}(t)=100-0.05t,
\qquad t\ \text{以秒计},
\label{eq:q3-capacity-time}
\end{equation}
可行性要求
\begin{equation}
m_h+m_c(n)\le100-0.05t,
\qquad m_h=60\ \mathrm{kg}.
\label{eq:q3-weight-constraint}
\end{equation}
因此仅由负重能力决定的最长站立时间为
\begin{equation}
t_W(n)
=\frac{100-m_h-m_c(n)}{0.05}
=20\left[40-m_c(n)\right].
\label{eq:q3-weight-time}
\end{equation}
在$n=0,1,2,3$时，$t_W(n)$随新增涂层数线性下降。该约束反映增加隔热厚度虽然降低热损失，但同时增加服装质量并缩短可承受站立时间。

\subsection{厚度变化后的传热正演模型}

对任意候选$n$，问题一的三层控制方程、人体能量平衡和两侧对流边界保持不变，仅将最外层区域厚度改为$d_3(n)$：
\begin{equation}
\left\{
\begin{aligned}
\rho_1c_1T_{1,t}&=k_1T_{1,xx},\\
\rho_2c_2T_{2,t}&=k_2T_{2,xx}+\rho_2q_{\rm pcm}(T_2),\\
\rho_3c_3T_{3,t}&=k_3T_{3,xx}.
\end{aligned}
\right.
\label{eq:q3-governing}
\end{equation}
人体侧仍取静止站立条件
\begin{equation}
-k_1T_{1,x}(0,t)=3.0\,[T_h(t)-T_1(0,t)],
\end{equation}
\begin{equation}
C_h\frac{\mathrm dT_h}{\mathrm dt}
=-3.0A_b[T_h(t)-T_1(0,t)],
\end{equation}
外侧仍采用问题一的自然对流边界
\begin{equation}
-k_3T_{3,x}(L(n),t)
=2.38\,[T_3(L(n),t)+40]^{1.25}.
\end{equation}
两个材料界面继续满足温度和热流连续，初始条件仍为人体及三层衣料统一$37\,^{\circ}\mathrm C$。

记方案$n$下由上述正演模型得到的人体温度首次达到$15\,^{\circ}\mathrm C$的时间为
\begin{equation}
t_T(n)=\inf\{t>0:T_h^{(n)}(t)\le15\,^{\circ}\mathrm C\}.
\label{eq:q3-thermal-time}
\end{equation}
$t_T(n)$表示由低温传热决定的热安全时间。增加隔热层厚度会同时改变最外层导热热阻和热容量，因此其具体增益由问题一的瞬态模型直接计算，不预先用稳态热阻公式替代。

\subsection{热安全--负重耦合的离散优化模型}

实验者能够实际持续站立的时间必须同时满足热安全条件和负重条件，因此定义方案$n$的有效持续时间。问题三最终转化为在费用约束下最大化有效持续时间的离散优化问题。

\subsubsection{模型汇总}

{\boldmath
\begin{equation}
\max_{n}\quad t_{\rm eff}(n)
\label{eq:q3-summary-obj}
\end{equation}}

{\boldmath
\begin{equation}
\mathrm{s.t.}\quad
\left\{
\begin{aligned}
&d_3(n)=0.3(1+n)\ \mathrm{mm},\\
&n\in\{0,1,2,3\},\\
&C_0 + n\Delta C \le 1.5 C_0,\\
&m_c(n) = A(\rho_1 d_1 + \rho_2 d_2 + \rho_3 d_3(n)),\\
&t_W(n) = 20\bigl[40 - m_c(n)\bigr],\\
&t_T(n) = \inf\{t>0: T_h^{(n)}(t) \le 15\,^\circ\mathrm{C}\},\\
&t_{\rm eff}(n) = \min\{t_T(n), t_W(n)\},\\
&m_h + m_c(n) \le 100 - 0.05 t_{\rm eff}(n).
\end{aligned}
\right.
\label{eq:q3-summary-con}
\end{equation}}

从模型结构上看，$t_T(n)$随隔热厚度增加而上升，而$t_W(n)$随$n$严格下降，因此最优方案通常出现在两类限制接近的位置或费用允许的边界处。由于候选方案只有四个，采用完全枚举即可获得严格的全局最优解：依次令$n=0,1,2,3$调用问题一正演求解器，计算$t_T(n)$、$t_W(n)$与$t_{\rm eff}(n)$，最后选择
\begin{equation}
n^*=\arg\max_{n\in\{0,1,2,3\}}t_{\rm eff}(n),
\end{equation}
对应的最优最外层厚度和总厚度分别为
\begin{equation}
d_3^*=0.3(1+n^*)\ \mathrm{mm},
\qquad
L^*=1.4+0.3n^*\ \mathrm{mm}.
\end{equation}
该趋势只作为机理判断，最终选择仍以四个候选方案的完整瞬态计算结果为准。

\subsection{模型求解算法}

本问优化问题的搜索空间仅包含四个离散候选方案，因此直接采用完全枚举法求解。算法核心流程见算法\ref{alg:q3-enumeration}。

\begin{algorithm}[H]
\SetAlgoLined
\setstretch{0.92}
\SetAlFnt{\normalsize}
\SetArgSty{textnormal}
\SetKwInput{KwIn}{Input}
\SetKwInput{KwOut}{Output}
\SetKwInput{KwInit}{Initialize}
\SetKwFor{For}{for}{do}{end for}
\SetKwIF{If}{ElseIf}{Else}{if}{then}{else if}{else}{end if}
\SetKw{KwRet}{Return}
\SetKw{KwCont}{continue}
\caption{双阈值离散厚度枚举算法}
\label{alg:q3-enumeration}
\KwIn{DSC扫描速率$\beta_{\rm DSC}$、放热曲线$(T,p(T))$及问题一校准后的传热模型参数}
\KwInit{候选集合$\mathcal N=\{0,1,2,3\}$，时间步$\Delta t=0.25\ \mathrm{s}$，模态数$M=40$}
\KwOut{各$n\in\mathcal N$的$t_{15}(n),t_{10}(n),t_W(n),t_{\rm eff,15}(n),t_{\rm eff,10}(n)$及最优层数$n_{15}^*,n_{10}^*$}

对每个$n\in\mathcal N$计算$d_3(n)=0.3(1+n)\times10^{-3}\ \mathrm{m}$\;
计算服装质量$m_c(n)=A(\rho_1d_1+\rho_2d_2+\rho_3d_3(n))$\;
计算负重限制时间$t_W(n)=20[40-m_c(n)]$\;
\For{each $n\in\mathcal N$}{
  初始化温度场为$37\,^\circ\mathrm{C}$，不可逆相变进度$\xi\leftarrow0$\;
  \While{$T_1(0,t)>10\,^\circ\mathrm{C}$ and $t<3600\ \mathrm{s}$}{
    更新外侧自然对流边界$-k_3T_{3,x}=2.38[T_3(L,t)+40]^{1.25}$\;
    由DSC数据与当前$\xi$计算功能层有效比热$c_{2,\rm eff}$\;
    推进一个时间步的模态解，并更新模态系数、平均温度和相变进度\;
    $t\leftarrow t+\Delta t$\;
  }
  插值得到$T_1(0,t)$首次达到$15\,^\circ\mathrm{C}$与$10\,^\circ\mathrm{C}$的时刻$t_{15}(n),t_{10}(n)$\;
  计算$t_{\rm eff,\theta}(n)=\min\{t_\theta(n),t_W(n)\}$，$\theta\in\{15,10\}$\;
}
对$\theta\in\{15,10\}$选择$n_\theta^*=\arg\max_{n\in\mathcal N}t_{\rm eff,\theta}(n)$\;
\KwRet{$n_{15}^*,n_{10}^*,d_3(n_\theta^*),t_{\rm eff,15}^*,t_{\rm eff,10}^*$及候选方案结果}\;
\end{algorithm}


\subsection{求解结果与分析}

本文对三种DSC扫描速率（$v=2, 5, 10\ \mathrm{K/min}$）分别计算四个候选方案，得到完整结果如表\ref{tab:q3-candidates}所示。最优方案汇总见表\ref{tab:q3-optimum}。

\begin{table}[!htbp]
\centering
\caption{问题三各候选方案计算结果}
\label{tab:q3-candidates}
\resizebox{\textwidth}{!}{%
\begin{tabular}{cccccccccc}
\toprule
\rowcolor{CumcmLightBlue}
$v$ (K/min) & $n$ & $d_3$ (mm) & $L$ (mm) & $m_c$ (kg) & $C$ (元) & $t_T$ (s) & $t_W$ (s) & $t_{\mathrm{eff}}$ (s) & 限制因素 \\
\midrule
\multirow{4}{*}{2}
& 0 & 0.3 & 1.4 & 0.754 & 301.67 & 520.90 & 784.92 & 520.90 & 热安全 \\
& 1 & 0.6 & 1.7 & 0.903 & 346.28 & 578.29 & 781.94 & 578.29 & 热安全 \\
& 2 & 0.9 & 2.0 & 1.052 & 390.89 & 638.33 & 778.97 & 638.33 & 热安全 \\
& 3 & 1.2 & 2.3 & 1.201 & 435.50 & 701.20 & 775.99 & 701.20 & 热安全 \\
\midrule
\multirow{4}{*}{5}
& 0 & 0.3 & 1.4 & 0.754 & 301.67 & 273.91 & 784.92 & 273.91 & 热安全 \\
& 1 & 0.6 & 1.7 & 0.903 & 346.28 & 318.76 & 781.94 & 318.76 & 热安全 \\
& 2 & 0.9 & 2.0 & 1.052 & 390.89 & 366.09 & 778.97 & 366.09 & 热安全 \\
& 3 & 1.2 & 2.3 & 1.201 & 435.50 & 416.04 & 775.99 & 416.04 & 热安全 \\
\midrule
\multirow{4}{*}{10}
& 0 & 0.3 & 1.4 & 0.754 & 301.67 & 191.64 & 784.92 & 191.64 & 热安全 \\
& 1 & 0.6 & 1.7 & 0.903 & 346.28 & 232.33 & 781.94 & 232.33 & 热安全 \\
& 2 & 0.9 & 2.0 & 1.052 & 390.89 & 275.43 & 778.97 & 275.43 & 热安全 \\
& 3 & 1.2 & 2.3 & 1.201 & 435.50 & 321.06 & 775.99 & 321.06 & 热安全 \\
\bottomrule
\end{tabular}%
}
\end{table}

\begin{table}[!htbp]
\centering
\caption{问题三不同扫描速率下的最优方案}
\label{tab:q3-optimum}
\begin{tabular}{ccccc}
\toprule
\rowcolor{CumcmLightBlue}
$v$ (K/min) & $n^*$ & $d_3^*$ (mm) & $t_{\mathrm{eff}}^*$ (s) & 限制因素 \\
\midrule
2 & 3 & 1.2 & 701.20 & 热安全 \\
5 & 3 & 1.2 & 416.04 & 热安全 \\
10 & 3 & 1.2 & 321.06 & 热安全 \\
\bottomrule
\end{tabular}
\end{table}

\FloatBarrier

从计算结果可以得到以下结论：

\noindent\textbf{1. 最优方案一致性}：在所有三种扫描速率下，最优方案均为$n^*=3$，即新增三层涂层，使最外层总厚度达到$1.2\ \mathrm{mm}$，总厚度达到$2.3\ \mathrm{mm}$。

\noindent\textbf{2. 限制因素分析}：在所有候选方案中，有效时间均由热安全约束决定，负重约束始终未被触发。这说明在当前参数条件下，增加隔热厚度对延长热安全时间的正效应，始终强于增重对缩短负重时间的负效应。

\noindent\textbf{3. 单调性验证}：热安全时间$t_T(n)$随新增层数$n$单调递增，而负重时间$t_W(n)$随$n$单调递减，验证了机理分析的正确性。

\noindent\textbf{4. 数值收敛性验证}：本问沿用问题一已验证的数值设置（时间步$0.25\ \mathrm{s}$，模态数$M=40$），对最优方案的收敛性检查表明：时间步减半改变结果不超过$0.004\%$，模态数增加改变结果不超过$0.3\%$，数值精度满足要求。